\documentclass[11pt]{article}

\usepackage{amsmath,amssymb}
\usepackage{geometry}
\usepackage{booktabs}
\usepackage{microtype}
\usepackage[hidelinks]{hyperref}

\geometry{margin=1in}

\title{A Spherical N-Series Area Law for Regular Refinement Families}
\author{C. Wayne Baker}
\date{June 03, 2026}

\begin{document}

\maketitle

\begin{abstract}
This paper presents a spherical extension of the N-Series geometric
cancellation framework.  The central claim is that regular, scale-consistent
spherical area approximations can exhibit an even-power error expansion under
refinement.  When this structure is present, fixed N-Series scale operators
cancel successive error terms and raise the observed convergence order.  Five
tests are reported.  Spherical caps, chordal spherical triangle patches,
latitude-longitude spherical grid cells, and fixed irregular spherical
polygon domains all exhibit the clean order ladder $2,4,6,8$ for the raw,
two-scale, three-scale, and four-scale approximations respectively.  A
perturbation test then identifies the regularity boundary: the raw
second-order law remains robust under smooth mesh perturbation, while
higher-order cancellation becomes sensitive to the loss of scale regularity.
The result is a geometric high-order area-refinement principle for regular
spherical refinement families, not a black-box method for arbitrary irregular
meshes.
\end{abstract}

\section{Introduction}

The N-Series framework begins from a simple observation: many geometric
approximations carry structured asymptotic error expansions.  When the
leading errors occur in even powers of the scale parameter, approximations at
several scales can be combined to cancel those leading terms.

In planar polygonal and perimeter problems, this produces the order ladder
\[
  2,\qquad 4,\qquad 6,\qquad 8,\ldots .
\]
The purpose of this paper is to show that the same structure appears in
spherical area approximation.

The tests considered here are deliberately elementary but geometrically
distinct:

\begin{itemize}
\item a spherical cap approximated by a geodesic boundary polygon;
\item a right spherical triangle approximated by a chordal triangulated patch;
\item a latitude-longitude spherical cell approximated by a chordal grid;
\item fixed irregular spherical polygon domains approximated by regular
      chordal fan refinements;
\item a perturbed latitude-longitude grid used to test robustness.
\end{itemize}

The first four tests confirm the clean N-Series order ladder for regular
refinement families, even when the domain boundary is irregular.  The fifth
test shows that high-order cancellation has a real regularity boundary when
the refinement mesh itself is perturbed.  This distinction is important: the
method should be stated as a high-order principle for regular,
scale-consistent refinement families, not as a universal correction for
arbitrary meshes.

The novelty of the present paper is not a new local quadrature formula for spherical area.  The contribution is the identification of a scale-structured cancellation law for spherical geometric approximations.  When the same spherical area approximation is evaluated on a regular refinement family at the scales \(N,2N,4N,\ldots\), the leading even-power error terms can be cancelled by fixed N-Series weights.  This separates the local construction of the approximation from the global cancellation mechanism, and explains why the same order ladder appears across caps, spherical triangle patches, latitude-longitude cells, and irregular spherical polygon domains.


\section{Main Contributions}

This paper makes four contributions.

First, it formulates a spherical N-Series area law for regular refinement
families:
\[
  A_N(D)
  =
  A(D)
  +
  \frac{C_2(D)}{N^2}
  +
  \frac{C_4(D)}{N^4}
  +
  \frac{C_6(D)}{N^6}
  +
  \cdots .
\]
The statement is conjectural in general, but is strongly supported by the
tests reported below.

Second, it shows that the fixed N-Series operators produce the expected
successive order improvements in several spherical settings:
\[
  O(N^{-2})
  \quad\longrightarrow\quad
  O(N^{-4})
  \quad\longrightarrow\quad
  O(N^{-6})
  \quad\longrightarrow\quad
  O(N^{-8}).
\]

Third, it separates boundary irregularity from mesh irregularity.  Fixed
irregular spherical polygon domains preserve the full $2,4,6,8$ ladder when
the refinement is regular and scale-consistent.

Fourth, it identifies a limitation: perturbing the refinement mesh preserves
the raw second-order law but can degrade higher-order cancellation.  Thus
regularity is a structural condition, not a cosmetic assumption.

\section{N-Series Scale Operators}

Let $D$ be a spherical domain on the unit sphere.  Let $A(D)$ denote its
exact spherical area, and let $A_N(D)$ denote a geometric area approximation
at scale $N$.  The spherical N-Series model is that, for suitable regular
refinement families,
\[
  A_N(D)
  =
  A(D)
  +
  \frac{C_2(D)}{N^2}
  +
  \frac{C_4(D)}{N^4}
  +
  \frac{C_6(D)}{N^6}
  +
  \cdots .
\]
The coefficients depend on the domain and on the approximation family, but
the powers are even.

The two-scale N-Series operator is
\[
  \widehat{A}^{(2)}_N(D)
  =
  A_{2N}(D)
  +
  \frac{A_{2N}(D)-A_N(D)}{3}.
\]
If the expansion above holds, then
\[
  \widehat{A}^{(2)}_N(D)
  =
  A(D)+O(N^{-4}).
\]

The three-scale operator is
\[
  \widehat{A}^{(3)}_N(D)
  =
  \frac{1}{45}A_N(D)
  -
  \frac{4}{9}A_{2N}(D)
  +
  \frac{64}{45}A_{4N}(D),
\]
and gives
\[
  \widehat{A}^{(3)}_N(D)
  =
  A(D)+O(N^{-6}).
\]

The four-scale operator is
\[
  \widehat{A}^{(4)}_N(D)
  =
  -\frac{1}{2835}A_N(D)
  +
  \frac{4}{135}A_{2N}(D)
  -
  \frac{64}{135}A_{4N}(D)
  +
  \frac{4096}{2835}A_{8N}(D),
\]
and gives
\[
  \widehat{A}^{(4)}_N(D)
  =
  A(D)+O(N^{-8}).
\]

Thus the method is not a new local formula for spherical area.  It is a
scale-cancellation operator acting on a family of geometric approximations.


\paragraph{Relation to classical extrapolation.}
The displayed weights are the Richardson--Romberg extrapolation weights for
scale doubling and an even-power error expansion.  Their numerical arithmetic
is classical: each weight set sums to one and annihilates the relevant powers
$N^{-2},N^{-4},\ldots$.  The contribution studied here is geometric rather
than algorithmic.  It is the evidence that several regular spherical area
families possess the required even-power structure, together with the observed
loss of a clean high-order ladder when scale coherence is disturbed.

\section{Spherical N-Series Area Law}

The numerical evidence motivates the following statement.

\medskip

\noindent
\textbf{Conjecture.}
Let $D$ be a spherical domain and let $A_N(D)$ be a regular,
scale-consistent geometric area approximation at scale $N$.  If the
refinement family preserves the same geometric structure across the scales
$N$, $2N$, $4N$, and $8N$, then the area error may admit an expansion of the
form
\[
  A_N(D)-A(D)
  =
  \frac{C_2(D)}{N^2}
  +
  \frac{C_4(D)}{N^4}
  +
  \frac{C_6(D)}{N^6}
  +
  \cdots .
\]
Under this even-power expansion, the N-Series operators produce the
successive order improvements
\[
  O(N^{-2})
  \quad\longrightarrow\quad
  O(N^{-4})
  \quad\longrightarrow\quad
  O(N^{-6})
  \quad\longrightarrow\quad
  O(N^{-8}).
\]

\medskip

The word regular is essential.  The perturbation study below shows that raw
second-order convergence can survive mesh perturbation, but higher-order
cancellation is sensitive to the loss of scale regularity.

\section{Evidence I: Spherical Cap Area}

For a spherical cap on the unit sphere with polar angle $\alpha$, the exact
area is
\[
  A_{\mathrm{cap}}
  =
  2\pi(1-\cos\alpha).
\]
The cap boundary was approximated by a geodesic polygon with $N$ equally
spaced boundary vertices.  The spherical polygon area was computed from
spherical angle excess.

The final observed orders at $N=256$ are shown in
Table~\ref{tab:cap-orders}.

\begin{table}[ht]
\centering
\small
\begin{tabular}{ccccc}
\toprule
$\alpha$ & Raw order & Two-scale order & Three-scale order & Four-scale order \\
\midrule
$15^\circ$ & $1.99989583$ & $3.99999065$ & $6.00093925$ & $7.99993479$ \\
$30^\circ$ & $1.99996740$ & $4.00041698$ & $6.00000053$ & $7.95025527$ \\
$60^\circ$ & $2.00016296$ & $4.00014936$ & $6.00003853$ & $7.99880645$ \\
$75^\circ$ & $2.00023455$ & $4.00028432$ & $6.00028420$ & $8.00026989$ \\
\bottomrule
\end{tabular}
\caption{Observed orders for spherical cap area at $N=256$.}
\label{tab:cap-orders}
\end{table}

The cap test confirms the N-Series ladder in the simplest spherical setting:
a rotationally symmetric domain with an exact closed-form area.


\paragraph{Why even powers arise for the cap.}
For the cap test, the even-power structure can be justified directly.  Let
$h=2\pi/N$, let $c=\cos\alpha$, and let $s=\sin\alpha$.  Joining the north
pole to two adjacent polygon vertices produces a spherical triangle whose
oriented area is
\[
  E_\alpha(h)
  =
  2\arctan\!\left(
    \frac{s^2\sin h}
         {1+2c+c^2+s^2\cos h}
  \right)
\]
for sufficiently small positive $h$.  The numerator is odd in $h$, the
denominator is even and positive near $h=0$, and therefore $E_\alpha(h)$ is
an odd analytic function:
\[
  E_\alpha(h)=a_1h+a_3h^3+a_5h^5+\cdots .
\]
The geodesic polygon area is $A_N=N E_\alpha(2\pi/N)$.  Since
$N=2\pi/h$, it follows that
\[
  A_N
  =2\pi a_1
   +\frac{(2\pi)^3a_3}{N^2}
   +\frac{(2\pi)^5a_5}{N^4}
   +\cdots .
\]
The constant term is the cap area $2\pi(1-\cos\alpha)$, so the cap family has
an even-power error expansion.  This argument is specific to the rotationally
symmetric cap construction; the broader spherical law remains conjectural.

\paragraph{Precision note.}
The $30^\circ$ four-scale order $7.95025527$ is slightly below the values near
$8$ seen in the other cap cases.  At this scale the four-scale residual is
very small, so a one-step observed-order quotient is sensitive to working
precision and near-cancellation.  The value is retained as measured rather
than adjusted.

\section{Evidence II: Chordal Spherical Triangle Area}

The second test uses a right spherical triangle on the unit sphere with
vertices
\[
  A=(0,0,1),
  \qquad
  B=(1,0,0),
  \qquad
  C=(\cos\beta,\sin\beta,0).
\]
For this triangle, the exact spherical area is
\[
  A_{\mathrm{exact}}=\beta,
\]
where $\beta$ is measured in radians.

The approximation uses a normalized barycentric grid.  Each small reference
triangle is mapped to the sphere and treated as a Euclidean chordal triangle
in $\mathbb{R}^3$.  The raw area $A_N$ is the sum of the chordal triangle
areas.

At $N=64$, the observed orders are shown in
Table~\ref{tab:triangle-orders}.

\begin{table}[ht]
\centering
\small
\begin{tabular}{ccccc}
\toprule
$\beta$ & Raw order & Two-scale order & Three-scale order & Four-scale order \\
\midrule
$15^\circ$ & $1.99942072$ & $3.99915618$ & $5.99886117$ & $7.99716985$ \\
$30^\circ$ & $1.99938698$ & $3.99908575$ & $5.99880136$ & $7.99761444$ \\
$60^\circ$ & $1.99923969$ & $3.99881849$ & $5.99858915$ & $7.99835976$ \\
$90^\circ$ & $1.99885313$ & $3.99827889$ & $5.99785193$ & $7.99808456$ \\
\bottomrule
\end{tabular}
\caption{Observed orders for chordal spherical triangle area at $N=64$.}
\label{tab:triangle-orders}
\end{table}

This confirms that the order ladder is not restricted to boundary polygons.
It also appears in a triangulated spherical patch.

\section{Evidence III: Latitude-Longitude Cell Area}

For a latitude-longitude cell on the unit sphere,
\[
  \lambda_0 \leq \lambda \leq \lambda_1,
  \qquad
  \phi_0 \leq \phi \leq \phi_1,
\]
the exact area is
\[
  A_{\mathrm{cell}}
  =
  (\lambda_1-\lambda_0)
  \bigl(\sin\phi_1-\sin\phi_0\bigr).
\]
The approximation uses a uniform latitude-longitude grid.  Each parameter
cell is split into two chordal triangles in $\mathbb{R}^3$, and the raw area
is the sum of the Euclidean triangle areas.

At $N=64$, the observed orders are shown in
Table~\ref{tab:latlon-orders}.

\begin{table}[ht]
\centering
\small
\begin{tabular}{lcccc}
\toprule
Cell & Raw order & Two-scale order & Three-scale order & Four-scale order \\
\midrule
equatorial $30^\circ\times30^\circ$ & $1.99996836$ & $3.99999637$ & $5.99994223$ & $7.99911861$ \\
midlatitude $30^\circ\times30^\circ$ & $1.99998486$ & $3.99996146$ & $5.99998582$ & $8.00085300$ \\
wide equatorial $60^\circ\times30^\circ$ & $1.99992674$ & $3.99992210$ & $6.00000864$ & $8.00036842$ \\
northern $15^\circ\times60^\circ$ & $1.99994200$ & $3.99996772$ & $5.99996443$ & $8.00007636$ \\
\bottomrule
\end{tabular}
\caption{Observed orders for latitude-longitude spherical cell area at $N=64$.}
\label{tab:latlon-orders}
\end{table}

This test is the most directly applied regular-grid example.  It connects
the N-Series area law to grid-cell geometry used in spherical mesh and
geodesy settings.

\section{Evidence IV: Irregular Spherical Polygon Domains}

The preceding tests used highly structured spherical domains: caps, right
spherical triangles, and latitude-longitude grid cells.  A natural question
is whether the N-Series ladder depends on the boundary being geometrically
simple.  To test this, fixed irregular convex spherical polygons were used.

For one six-vertex example, the polygon vertices were
\[
  (4^\circ,-8^\circ),\quad
  (18^\circ,5^\circ),\quad
  (15^\circ,28^\circ),\quad
  (2^\circ,42^\circ),\quad
  (-14^\circ,30^\circ),\quad
  (-10^\circ,4^\circ).
\]
Its exact spherical area, computed by spherical angle excess, was
\[
  A_{\mathrm{exact}}
  =
  0.325062777898240815426301649532.
\]
The approximation used a regular chordal fan refinement.  A root vertex was
chosen, the polygon was split into fan triangles, and each fan triangle was
approximated by a normalized barycentric chordal patch at scale $N$.

The observed orders for this fixed polygon are shown in
Table~\ref{tab:single-irregular-polygon-orders}.

\begin{table}[ht]
\centering
\small
\begin{tabular}{ccccc}
\toprule
$N$ & Raw order & Two-scale order & Three-scale order & Four-scale order \\
\midrule
$16$ & $1.99929248$ & $3.99269841$ & $5.99184034$ & $8.01301620$ \\
$32$ & $1.99982228$ & $3.99817005$ & $5.99798552$ & $8.00310155$ \\
$64$ & $1.99995552$ & $3.99954224$ & $5.99949790$ & $8.00076594$ \\
\bottomrule
\end{tabular}
\caption{Observed convergence orders for a fixed irregular spherical polygon
under regular chordal fan refinement.}
\label{tab:single-irregular-polygon-orders}
\end{table}

The full ladder is preserved.  At $N=64$, the observed orders are
\[
  1.99995552,\qquad
  3.99954224,\qquad
  5.99949790,\qquad
  8.00076594.
\]

The same test was then repeated across deterministic irregular polygon
families with $5$, $6$, and $8$ vertices and seeds $1$, $2$, and $3$.  The
aggregate final-level observed orders are shown in
Table~\ref{tab:batch-irregular-polygon-orders}.

\begin{table}[ht]
\centering
\small
\begin{tabular}{ccccc}
\toprule
Vertices & Raw order & Two-scale order & Three-scale order & Four-scale order \\
\midrule
$5$ & $1.99993053$ & $3.99981653$ & $5.99991832$ & $7.99955672$ \\
$6$ & $1.99993586$ & $3.99979158$ & $5.99991390$ & $8.00029966$ \\
$8$ & $2.00001032$ & $4.00140962$ & $5.99925770$ & $7.99912295$ \\
\bottomrule
\end{tabular}
\caption{Aggregate final-level observed orders for deterministic irregular
spherical polygon families under regular chordal fan refinement.}
\label{tab:batch-irregular-polygon-orders}
\end{table}

This confirms that the irregular-boundary result is not a single-polygon
accident.  Boundary irregularity alone does not destroy the N-Series ladder
when the refinement remains regular and scale-consistent.

\section{Evidence V: Mesh Perturbation and Regularity Boundary}

The final test perturbs the interior nodes of a latitude-longitude cell.  The
boundary nodes are fixed, while the interior grid points are displaced by a
smooth deterministic field.  The perturbation amplitude is measured as a
fraction of the local grid spacing.  Three seeds were tested at each
perturbation level.

The aggregate final-level mean orders are shown in
Table~\ref{tab:perturb-orders}.

\begin{table}[ht]
\centering
\small
\begin{tabular}{ccccc}
\toprule
Perturbation & Raw order & Two-scale order & Three-scale order & Four-scale order \\
\midrule
$0.00$ & $1.99996836$ & $3.99999637$ & $5.99994223$ & $7.99911861$ \\
$0.02$ & $1.99997529$ & $4.02004815$ & $4.53899093$ & $2.61899119$ \\
$0.05$ & $2.00001207$ & $3.44420303$ & $5.68235637$ & $2.77050496$ \\
$0.10$ & $2.00014360$ & $3.89815337$ & $8.94059560$ & $3.59334301$ \\
$0.20$ & $2.00066838$ & $3.90830154$ & $6.15480240$ & $4.36890149$ \\
\bottomrule
\end{tabular}
\caption{Mean final-level observed orders for perturbed latitude-longitude
cell area.}
\label{tab:perturb-orders}
\end{table}

The perturbation table must be interpreted as a diagnostic rather than as a
list of stable asymptotic orders.  Two conclusions are robust in the tested
window.  First, the raw approximation remains close to second order.  Second,
the two-scale operator remains broadly effective, although it is less uniform
than on the unperturbed grid.

The higher columns do not support a uniform sixth- or eighth-order claim under
perturbation.  In particular, the three-scale value $8.94059560$ exceeds the
nominal order of that operator.  Such a value can arise from transient
coefficient cancellation, a sign change in the residual, proximity to a
numerical floor, or loss of a common asymptotic expansion across scales.  It
is therefore not evidence of genuine order greater than six.  More generally,
the nonmonotone higher-order values show that the single-step statistic
\[
  p_N=\frac{\log(e_N/e_{2N})}{\log 2}
\]
is not measuring a stable asymptotic regime for those perturbed cases.

The conservative conclusion is that smooth perturbation preserves the raw
second-order behavior in these experiments and often leaves the two-scale
correction useful, but the data do not establish a clean higher-order ladder
for the perturbed meshes.  Regularity and scale coherence are therefore part
of the hypothesis of the spherical N-Series area law.

\section{Boundary Irregularity Versus Mesh Irregularity}

The irregular-polygon and perturbation tests separate two different
phenomena.

Boundary irregularity means that the spherical domain itself is not a cap,
triangle, or latitude-longitude rectangle.  The batch irregular-polygon test
shows that this does not destroy the N-Series ladder when the interior
refinement is still regular.

Mesh irregularity means that the refinement family itself loses coherent
scale structure.  The perturbation test shows that this can degrade the
higher-order operators, especially the four-scale operator.

Thus the decisive condition is not that the boundary must be simple.  The
decisive condition is that the refinement family must preserve scale
regularity.

\paragraph{Boundary stress sweep and failure mode.}
A small stress sweep was also performed by varying the number of polygon
vertices, the angular radius of the domain, and the boundary jitter.  For
well-posed irregular polygons, the full N-Series ladder persisted.  For
example, five-vertex cases at radii $15^\circ$ and $30^\circ$, with jitter
levels $0.2$ and $0.5$, retained observed orders close to
\[
  2,\qquad 4,\qquad 6,\qquad 8.
\]
The same was true for the eight-vertex case at jitter $0.2$.

However, the eight-vertex cases with jitter $0.5$ did not converge to the
same area under the chordal fan approximation.  The raw error remained nearly
constant under refinement, and the observed orders collapsed to approximately
zero.  This indicates a geometric pathology in the polygon/fan pairing, such
as nonconvexity, orientation failure, self-overlap, or convergence to a
different spherical region.  This is not a failure of N-Series cancellation
on a valid refinement family; it is a failure of the approximation family to
target the intended domain.

The stress sweep therefore sharpens the regularity condition: boundary
irregularity is allowed, but the domain and refinement must remain
geometrically well posed.

\section{Limitations}

The results in this paper are numerical and should be interpreted as
evidence for a geometric area law, not as a proof for arbitrary spherical
domains.

The following limitations are important:

\begin{itemize}
\item all tests use convex or well-behaved spherical domains;
\item the exact areas are available from closed-form formulas or spherical
      angle excess;
\item the successful high-order cases use regular, scale-consistent
      refinement families;
\item the perturbation study shows that high-order cancellation may degrade
      when mesh regularity is lost;
\item no claim is made for arbitrary nonconvex, self-intersecting, adaptive,
      or highly irregular spherical meshes.
\end{itemize}

These limitations are part of the claim.  The method is a high-order
geometric cancellation principle for regular refinement families, not a
universal mesh repair algorithm.

\section{Relation to the General N-Series Framework}

The spherical area law extends the N-Series framework from planar perimeter
and chord approximation to curved-surface area.  The operator mechanism is
unchanged: one computes the same geometric quantity at several scales and
combines those values with fixed weights so that the leading error modes
vanish.

The result is significant because it shows that N-Series cancellation is not
limited to approximating lengths or constants such as $\pi$.  It also applies
to area approximation on curved surfaces when the refinement family preserves
sufficient structure across scales.

\section{Conclusion}

The conservative conclusion is that spherical N-Series cancellation provides
a geometric high-order area-refinement method for regular spherical
refinement families.  Boundary irregularity is allowed, provided the
refinement remains regular and scale-coherent.  Thus the relevant hypothesis
is not smoothness of the boundary alone, but coherence of the entire
approximation family: the domain, orientation, subdivision rule, and target
area must remain fixed across refinement.

The method should not be claimed as a black-box correction for arbitrary
irregular meshes.  Within its regularity class, however, the evidence
strongly supports a spherical extension of the N-Series geometric
cancellation framework.

\appendix

\section{Reproducibility Notes}

The numerical results in this paper were generated by standalone Python
scripts.  Each script computes exact spherical reference areas, constructs a
geometric approximation at several refinement scales, applies the N-Series
operators, and reports observed convergence orders.

\subsection{Spherical Cap Test}

The spherical cap results were generated by
\[
  \texttt{spherical\_area\_nseries\_cap\_test.py}.
\]
For a cap of polar angle $\alpha$ on the unit sphere, the exact area is
\[
  A_{\mathrm{cap}}
  =
  2\pi(1-\cos\alpha).
\]

\subsection{Chordal Spherical Triangle Test}

The spherical triangle results were generated by
\[
  \texttt{spherical\_triangle\_nseries\_area\_test.py}.
\]
The target is a right spherical triangle with exact area
\[
  A_{\mathrm{exact}}=\beta.
\]

\subsection{Latitude-Longitude Cell Test}

The regular latitude-longitude cell results were generated by
\[
  \texttt{spherical\_latlon\_cell\_nseries\_area\_test.py}.
\]
For a unit-sphere cell, the exact area is
\[
  A_{\mathrm{cell}}
  =
  (\lambda_1-\lambda_0)
  \bigl(\sin\phi_1-\sin\phi_0\bigr).
\]

\subsection{Irregular Spherical Polygon Tests}

The irregular spherical polygon results were generated by
\[
  \texttt{spherical\_irregular\_polygon\_nseries\_test.py}
\]
and
\[
  \texttt{spherical\_irregular\_polygon\_batch\_nseries\_test.py}.
\]
The boundary stress sweep was generated by
\[
  \texttt{spherical\_irregular\_polygon\_stress\_sweep\_nseries\_test.py}.
\]
The exact area is computed by spherical angle excess.  The approximation uses regular chordal fan refinement.  The stress-sweep script varies the number of vertices, angular radius, and boundary jitter in order to distinguish well-posed irregular domains from pathological polygon/fan pairings.

\subsection{Perturbed Latitude-Longitude Cell Test}

The perturbation robustness results were generated by
\[
  \texttt{spherical\_perturbed\_latlon\_nseries\_test.py}.
\]
This script perturbs the interior grid nodes by a smooth deterministic field while keeping the boundary fixed.

\subsection{Observed Orders}

The reported convergence orders are computed from successive errors by
\[
  p_N
  =
  \frac{\log(e_N/e_{2N})}{\log 2},
\]
where $e_N$ is the absolute error at scale $N$.

\begin{thebibliography}{9}

\bibitem{coxeter}
H. S. M. Coxeter,
\emph{Introduction to Geometry}.
Wiley, 1969.

\bibitem{todhunter}
I. Todhunter,
\emph{Spherical Trigonometry}.
Macmillan, 1886.

\bibitem{richardson}
L. F. Richardson,
``The approximate arithmetical solution by finite differences of physical
problems,''
\emph{Philosophical Transactions of the Royal Society A}, 210, 1911,
pp.\ 307--357.

\bibitem{williamson}
D. L. Williamson,
``The Evolution of Dynamical Cores for Global Atmospheric Models,''
\emph{Journal of the Meteorological Society of Japan}, 85B, 2007,
pp.\ 241--269.

\end{thebibliography}

\end{document}
